% Pursuing Stacks (cotes 134-2 to 134-8): archive cote 134-6, site page 74 (PDF page 75), typescript page 392.
% typescript. Draft compared: ps5.tex:5169-5247.
% AI-transcribed by Claude Opus 5.5 (claude-opus-5-5, adaptive thinking, effort high), Message Batches API, batch ps3 (msgbatch_01XdFmuVVBAcgc7EFxdtCqsj), 2026-11-12; not yet checked by eye.
\scanpage{134-6/74}
\mspage{392}

(28) \[ \underline{Z}_{\hand{\mathcal{A}}} \hand{\circledast} L \simeq L \quad . \]

Using the tautological duality relation (20) between categories of
cosheaves (with values in \hand{$\mathcal{M}$}) and categories of sheaves (with values in
\hand{$\mathcal{N}=\mathcal{M}^{o}$}), we deduce accordingly an associative and unitary operation
of \hand{$\mathcal{A}$}$_{ab}$ on any category of the type \hand{$\mathcal{A}_{\mathcal{N}}$}, namely \hand{$\mathcal{N}$}-valued sheaves on \hand{$\mathcal{A}$}.
This operation is most conveniently denoted by the \uline{Hom} symbol
(29) \[ (F, K) \hand{\longmapsto} \underline{Hom}(F,K) \ : \ \hand{\mathcal{A}}_{ab}{}^{o} \, x \, \hand{\mathcal{A}_{\mathcal{N}} \longrightarrow \mathcal{A}_{\mathcal{N}}} \quad , \]
it's explicit description in terms of K\hand{,} viewed as a functor
\[ K\colon \hand{\mathcal{A}}_{ab}{}^{o} \hand{\longrightarrow} \hand{\mathcal{N}} \]
is by
(30) \[ \underline{Hom}(F,K)\ (G) \ \hand{\unsure{\simeq}}\xout{\unsure{Hom}} \ K(G \otimes F\ ) \]
for K in \hand{$\mathcal{A}_{\mathcal{N}}$}, F and G in \hand{$\mathcal{A}$}$_{ab}$ . This may be written more suggestively as
(30') \[ Hom(G,\ \underline{Hom}(F,K)) \hand{\simeq} Hom(G\otimes F,\ K) \qquad , \]
similar to (22'), where we use the notation Hom (non underlined !) in
analogy to (23) for denoting K (H) , namely
(31) \[ Hom(H,K) \overset{def}{=} K(H) \qquad . \]
If confusion is feared, we may \deleted{write the}\inserted{put a} subscript \uline{Z} in all Hom's and
\uline{Hom}'s just introduced, as well as in the internal and external tensor
product operations. \hand{$\otimes$, $*$, $\circledast$}
\inserted{\uline{Comments.}}
\hand{(}Take for instance
\[ \hand{\mathcal{N}} = (Ab) \ , \quad hence \ \hand{\mathcal{A}_{\mathcal{N}} \simeq \deleted{\illegible{1}} \mathcal{A}_{ab}} \]
by prop.1 (compare (14)). By this equivalence, the Hom in (31) is\hand{,} just
the usual Hom set corresponding to the category structure of \hand{$\mathcal{A}$}$_{ab}$ \inserted{(the Hom} endowed
moreover with the structure of abelian group, coming from the fact that
\hand{$\mathcal{A}$}$_{ab}$ is\xout{\illegible{1}} an additive category), whereas formula (30') shows that the
``external'' \uline{Hom} in this case is nothing but the \xout{inter} usual internal \uline{Hom}
as contemplated in (24). This I hope will convince the reader of the
adequacy of the notation used\over{.}{,} \hand{(}in case of exterior \xout{\illegible{1}} operation of
abelian sheaves on \hand{$\mathcal{N}$}-valued sheaves\hand{)}, and \inserted{of} the convention (31). We would
like to give a similar justification for the notations (22) (23) used in
connection with operation of abelian sheaves on \hand{$\mathcal{M}$}-valued \uline{co}sheaves,
by interpreting \hand{this} as \xout{an}\inserted[typed]{the} internal \xout{operation} tensor product operation in
\hand{$\mathcal{A}$}$_{ab}$ , for suitable choice of \hand{$\mathcal{N}$}. This I am affraid cannot be done for

\margin{74}
